% =====================================================================
\section{\texorpdfstring{\IIC{}}{I2C} controller}
\label{sec:i2c}
% =====================================================================

The \IIC{} controller operates as a bus master (7-bit addressing, 100 and
400\,kHz) and offers two levels of use: six SFRs for full control of a
transaction, and five high-level instructions (\texttt{I2CW},
\texttt{I2CR}, \texttt{I2CRR}, and the multi-byte \texttt{I2CRRN},
\texttt{I2CWN}) that cover the common cases in a single line of code.

\subsection{Rationale for the register set}

Three registers alone (slave address, byte to send, received byte) are enough
to write or read a single byte, but they do not allow:

\begin{itemize}
  \item reading a sensor register, which requires writing the register number
        followed by a \textbf{repeated START} without a STOP;
  \item multi-byte transactions (e.g.\ an EEPROM page write, reading 6 bytes
        from an accelerometer);
  \item error detection: missing ACK (device not responding), arbitration
        lost, a stuck bus;
  \item setting the SCL frequency and checking whether a transfer has finished.
\end{itemize}

The three registers are therefore kept (I2CADR, I2CTXD, I2CRXD), and a
control, a status and a clock divider register are added.

\subsection{Registers}

\needspace{14\baselineskip}
\paragraph{I2CCON (\hex{10})}

\begin{center}
\small
\begin{tabularx}{\linewidth}{@{}l l X@{}}
  \toprule
  \textbf{Bit} & \textbf{Name} & \textbf{Meaning} \\
  \midrule
  7 & EN    & Controller enable \\
  6 & IE    & Interrupt when an operation completes (DONE) \\
  5 & START & Write 1: START (or repeated START if the bus is already owned)
              followed by the address byte \{I2CADR, RW\}; cleared by hardware \\
  4 & STOP  & Write 1: STOP condition; cleared by hardware \\
  3 & RD    & Write 1: receive one byte into I2CRXD; cleared by hardware \\
  2 & ACKS  & Acknowledge sent after a received byte: 0 = ACK, 1 = NACK (last byte) \\
  1 & RW    & Direction in the address phase: 0 = write, 1 = read \\
  0 & --    & Reserved \\
  \bottomrule
\end{tabularx}
\end{center}

\needspace{14\baselineskip}
\paragraph{I2CSTA (\hex{11})}

\begin{center}
\small
\begin{tabularx}{\linewidth}{@{}l l X@{}}
  \toprule
  \textbf{Bit} & \textbf{Name} & \textbf{Meaning} \\
  \midrule
  0 & TIP   & Byte transfer in progress \\
  1 & DONE  & Operation complete (R/W1C); interrupt source \\
  2 & NACK  & The slave did not acknowledge the last byte \\
  3 & ARBL  & Arbitration lost (R/W1C) \\
  4 & BB    & Bus busy (between START and STOP) \\
  5 & RXRDY & New byte in I2CRXD; cleared by reading I2CRXD \\
  6 & TOUT  & Timeout expired (R/W1C) \\
  7 & --    & Reserved \\
  \bottomrule
\end{tabularx}
\end{center}

Writing I2CTXD while BB = 1 starts transmission of that byte. The SCL
frequency is
\begin{equation}
  f_{\mathrm{SCL}} = \frac{f_{\mathrm{CLK}}}{4 \cdot (\mathrm{I2CBRG} + 1)}
  \label{eq:scl}
\end{equation}

Example for $f_{\mathrm{CLK}} = 16$\,MHz: I2CBRG = 39 gives 100\,kHz,
I2CBRG = 9 gives 400\,kHz.

\subsection{High-level instructions}
\label{sec:i2c-hl}

\begin{center}
\small
\begin{tabularx}{\linewidth}{@{}l X@{}}
  \toprule
  \textbf{Instruction} & \textbf{Bus sequence} \\
  \midrule
  \texttt{I2CW Ra, Rd}  & START, address(Ra)+W, byte Rd, STOP \\
  \texttt{I2CR Rd, Ra}  & START, address(Ra)+R, byte $\rightarrow$ Rd (NACK), STOP \\
  \texttt{I2CRR Rd, Ra} & START, address(Ra)+W, byte Ra+1, repeated START,
                          address(Ra)+R, byte $\rightarrow$ Rd (NACK), STOP \\
  \texttt{I2CRRN [pair], Ra, \#n} & START, address(Ra)+W, byte Ra+1, repeated START,
                          address(Ra)+R, n bytes $\rightarrow$ mem[pair\ldots pair+n$-$1]
                          (ACK after each byte, NACK after the last), STOP \\
  \texttt{I2CWN [pair], Ra, \#n} & START, address(Ra)+W, n bytes from
                          mem[pair\ldots pair+n$-$1], STOP \\
  \bottomrule
\end{tabularx}
\end{center}

The multi-byte instructions transfer data directly between the \IIC{}
controller and data memory; the bytes are stored in bus order at increasing
addresses. To write device registers with \texttt{I2CWN}, the first byte of
the buffer holds the register number. Because the processor stalls for the
whole transaction (about 9 SCL periods per byte: roughly 0.8\,ms for 6 bytes
at 100\,kHz), longer transfers \SHOULD{} use 400\,kHz.

Example: reading 6 bytes (registers \hex{3B}--\hex{40}) from an
accelerometer at address \hex{68} into a buffer:

\begin{lstlisting}
    MOV    R2, #0x68        ; slave address
    MOV    R3, #0x3B        ; first register
    MOV    R4, #<buf        ; R5:R4 = buffer address
    MOV    R5, #>buf
    I2CRRN [R4], R2, #6     ; buf[0..5] <- registers 0x3B..0x40
    JC     i2c_error
\end{lstlisting}

\begin{figure}[htbp]
  \centering
  % ---------------------------------------------------------------------
%  I2C frame generated by I2CRR
%  Each field: x / width / driven-by-slave (1|0) / label
%  1 unit = 0.2 mm, y grows downwards.
% ---------------------------------------------------------------------
\begin{tikzpicture}[
    x=0.2mm, y=-0.2mm,
    conn/.style={draw=lineclr, line width=0.6pt},
    lbl/.style={font=\scriptsize, text=black!60},
    master/.style={draw=lineclr, line width=0.6pt, rounded corners=0.6mm},
    slave/.style={draw=accent, fill=accent!14, line width=0.9pt, rounded corners=0.6mm},
  ]
  \foreach \x/\w/\s/\t in {
      24/40/0/S,
      64/110/0/{address + W},
      174/40/1/A,
      214/70/0/Ra+1,
      284/40/1/A,
      324/40/0/Sr,
      364/110/0/{address + R},
      474/40/1/A,
      514/100/1/{data $\rightarrow$ Rd},
      614/40/0/N,
      654/40/0/P} {
    \ifnum\s=1
      \draw[slave] (\x,72) rectangle ++(\w,40);
    \else
      \draw[master] (\x,72) rectangle ++(\w,40);
    \fi
    \node[font=\scriptsize] at (\x+\w/2,92) {\t};
  }

  % --- phases ---------------------------------------------------------
  \draw[conn] (24,120) -- (24,128) -- (322,128) -- (322,120);
  \node[lbl, anchor=north] at (173,132) {write phase: register number};
  \draw[conn] (326,120) -- (326,128) -- (694,128) -- (694,120);
  \node[lbl, anchor=north] at (510,132) {read phase: one byte};

  % --- legend ---------------------------------------------------------
  \draw[slave]  (24,170)  rectangle ++(16,16);
  \node[font=\scriptsize, anchor=west] at (46,178)  {driven by slave};
  \draw[master] (180,170) rectangle ++(16,16);
  \node[font=\scriptsize, anchor=west] at (202,178) {driven by master (CISC-8)};

  \node[lbl, anchor=west] at (24,210)
    {S / P = START / STOP, Sr = repeated START, A = ACK, N = NACK};
  \node[lbl, anchor=west] at (24,230)
    {Missing ACK anywhere: STOP, I2CSTA.NACK = 1, flag C = 1};
\end{tikzpicture}

  \caption{\IIC{} frame generated by \texttt{I2CRR}: register number write, repeated START, single-byte read}
  \label{fig:i2cframe}
\end{figure}

The repeated START (Sr) keeps ownership of the bus between writing the
register number and reading the data, so no other master can intervene
(Figure~\ref{fig:i2cframe}).

Example: reading register \hex{75} from a sensor at address \hex{68}:

\begin{lstlisting}
    MOV   R2, #0x68     ; slave address
    MOV   R3, #0x75     ; register number
    I2CRR R0, R2        ; R0 <- data read
    JC    i2c_error     ; C = 1: NACK / timeout
\end{lstlisting}

\subsection{Requirements}

\begin{reqtable}
\req{I2C-001}{The controller \MUST{} operate as a master only, with 7-bit
  addressing, in accordance with NXP UM10204. Slave mode and 10-bit
  addressing are not supported.}
\req{I2C-002}{The controller \MUST{} support Standard mode (100\,kHz) and Fast
  mode (400\,kHz); the frequency is set by I2CBRG according to
  Equation~\eqref{eq:scl}.}
\req{I2C-003}{The controller \MUST{} generate START, repeated START and STOP
  conditions, and transmit and receive bytes with ACK/NACK handling.}
\req{I2C-004}{The controller \MUST{} support clock stretching by the slave by
  sampling the SCL line.}
\req{I2C-005}{The controller \MUST{} detect arbitration loss (SDA read back
  differs from SDA driven), release the bus and set ARBL.}
\req{I2C-006}{The controller \MUST{} abort the operation and set TOUT when SCL
  is held low longer than the timeout (synthesis parameter, default 25\,ms).}
\req{I2C-007}{SDA and SCL \MUST{} be driven as open-drain (low or
  high-impedance only); the pull-up resistors are external.}
\req{I2C-008}{The SDA and SCL inputs \MUST{} be synchronised to CLK (at least
  2 flip-flops) and \SHOULD{} have a glitch filter suppressing pulses shorter
  than 50\,ns.}
\req{I2C-009}{Setting DONE while IE = 1 \MUST{} raise the \IIC{} interrupt.}
\req{I2C-010}{The instructions \texttt{I2CW}, \texttt{I2CR}, \texttt{I2CRR},
  \texttt{I2CRRN} and \texttt{I2CWN} \MUST{} use the same controller and
  leave the final status in I2CSTA.}
\req{I2C-011}{For \texttt{I2CRRN} and \texttt{I2CWN} the controller \MUST{}
  read and write data memory directly, one byte per transferred \IIC{} byte
  (\reqref{ISA-041}).}
\end{reqtable}
